In the context of differentially private machine learning, \emph{privacy amplification by subsampling} exploits the randomness in the selection of minibatches to provide better privacy guarantees than if the minibatches were fixed.
This is typically analyzed in the setting where we add independent noise to each minibatch update.
A recent line of work has shown that adding correlated noise across updates using \emph{the matrix mechanism} can give better results than independent noise in practical settings.
However, correlated noise is not amenable to standard analyses of privacy amplification by subsampling.
Thus practitioners must choose between the benefits of correlated noise and the benefits of privacy amplification. In this work, we show how we can enjoy the benefits of both.

To be more precise, we consider $\dimension$ rounds\footnote{Note that we use $n$ to denote the number of rounds rather than dataset size, following the notation in \cite{denisovDPMF, MEMF, BandedMF}.}. 
Each round has a corresponding minibatch such that each example is included in the minibatch with probability $p$, and an example's inclusion in each minibatch is independent of other rounds and other examples.  
In each round, we compute the gradient of the model's loss on the examples in the minibatch, clip these per-example gradients to have 2-norm $\le 1$, sum these clipped gradients, add some noise, and finally use this noisy aggregate gradient to update the model parameters.
The noise is generated by the matrix mechanism, which is described by a lower-triangular, full-rank, non-negative matrix $\bfC \in \mathbb{R}_{\geq 0}^{\dimension \times \dimension}$.
The noise is independent across the coordinates of the gradients/model parameters, but correlated across the rounds. Specifically, for a given gradient coordinate/model parameter, the noise added across the $\dimension$ rounds is generated by $\bfC^{-1} \bfz$, where $\bfz \sim N(0, \sigma^2)^\dimension$ is $\dimension$ independent Gaussians.
Equivalently, the noise comes from $N(0,\Sigma)$ with $\Sigma = \sigma^2 \bfC^{-1} (\bfC^{-1})^T$.
%In the $j$th round, we receive the $j$th row of a matrix $\obs$, which is a function of the database $\dataset$, and release the $j$th row of $\bfC \bfx + \bfz$, where $\bfz \sim N(0, \sigma^2)^\dimension$. With amplification by sampling, each row of $\obs$ is a function of a subsample of $\dataset$, chosen such that each example is included with probability $p$, and an example's inclusion in each round is independent of other rounds. 
Adding independent noise across rounds corresponds to $\bfC=I$.
The final privacy guarantee depends on the number of rounds $\dimension$, the subsampling probability $p$, the noise multiplier $\sigma$, and, of course, the matrix $\bfC$.%\tnote{Is there a reason why the number of rounds is $n$? usually $n$ denotes the number of people?}\agnote{It's a holdover from the matrix-factorization papers, which started by operating in the single-pass streaming setting, i.e. rounds = dataset size. Happy to pivot to another term.}

\subsection{Prior Work}

The independent noise setting (i.e., $\bfC=I$, a.k.a.~DP-SGD) is extremely well-studied \citep{song2013stochastic,BST14,abadi2016deep,mironov2019r,steinke2022composition,koskela2020computing} and we have tight analyses.
In particular, the subsampling guarantee is essentially optimal in the regime where the final privacy guarantee is strong (e.g., $\varepsilon\le1$), which means that the noise multiplier $\sigma$ is sufficiently large.
However, in practice, the final privacy guarantee is often not so strong (e.g., $1<\varepsilon\le10$) and the noise multiplier $\sigma$ is not large enough to give good privacy amplification. This means we require large minibatches to achieve good performance.

To address the limitations of privacy amplification by subsampling, a recent line of work \citep{kairouz2021practical,denisovDPMF,MEMF,xu2023federated} has developed algorithms that add correlated noise across rounds. These algorithms build on work on continual observation \citep{dwork2010differential}.
Intuitively, the benefit of correlated noise is that the noise we add over multiple rounds partially cancels out.
This yields better results than independent noise in practical settings.

The downside of using the matrix mechanism is that it does not readily benefit from privacy amplification by subsampling. That is, the standard privacy guarantee does not depend on the subsampling probability $p$. And it is not obvious how to combine privacy amplification by subsampling with correlated noise.

\citet{BandedMF} analyzed a special case where the matrix $\bfC$ is ``banded'' so that $\bfC_{i,j}=0$ if $i<j+b$. Intuitively, this means that rounds that are more than $b$ steps apart are independent. This allows us to reduce to the usual privacy amplification analysis (with a loss of a factor of $b$).
However, in many cases we do not have such a banded structure.
Furthermore, the banded analysis is brittle -- a single nonzero entry far from the diagonal breaks this banded structure -- yet we do not expect that the true privacy guarantee of the algorithm is so brittle.

%Mention of DP-MF, continual observation works. DP-MF used to not be better than DP-SGD in low epsilon regime because it doesn't benefit from amplification. Banded paper closes the gap, but not completely.

%Comparison to amplification argument from banded paper. Idea: Banded paper has two main weaknesses. The first is that e.g. if we take the identity and make every entry below the diagonal a small positive value, the mechanism's privacy guarantees should be basically the same, but the banded paper's guarantees get substantially worse. In addition, there's more randomness if we use a smaller sampling probability $p$ over $n$ rounds (instead of $bp$ instead of $n/b$ rounds), i.e. more potential for privacy amplification, but the banded paper requires the less random sampling scheme for the analysis / does not get full benefit of privacy amplification. 

\subsection{Our contributions}

Our main contribution is a method for analyzing the matrix mechanism with privacy amplification by subsampling, which is applicable to general correlation matrices $\bfC$. 
We demonstrate that our method yields both asymptotic theoretical improvements and practical experimental improvements.
%We give a summary of our main challenges and contributions, as well as the structure of our paper here.

\mypar{Conditional composition (\cref{sec:conditioningtrick})} 
Our main technical tool is a version of the composition theorem that gracefully handles dependence between rounds. 
Standard composition theorems can handle dependence between rounds, but only by assuming a worst-case privacy guarantee for each round that holds regardless of what happened in prior rounds; this is too pessimistic.
Our conditional composition theorem (\cref{thm:conditioningtrick}) instead assumes a privacy guarantee for each round which holds with high probability over what happened in prior rounds. These high-probability privacy guarantees can be significantly better than the worst-case guarantees, which allows us to obtain better results for the matrix mechanism.  %challenge is that in the matrix mechanism, each round is dependent, whereas composition assumes independence in the non-adaptive case. Tools like adaptive composition can handle dependence between rounds, if the worst-case privacy guarantee of round $i$ conditioned on the previous rounds is bounded. However, for the matrix mechanism, in the worst-case the adversary may learn exactly which examples participated in rounds 1 to $i-1$ prior to seeing round $i$'s output. In contrast, if we only released the output of round $i$, the examples participating in these rounds fully benefits from privacy amplification. So, the worst-case privacy guarantee of round $i$ conditioned on the previous rounds will be much worse than the ``average-case,'' and thus adaptive composition will give weak privacy guarantees.
%In \cref{thm:conditioningtrick}, we give a \textit{conditional composition theorem}, a generalization of adaptive composition, to handle composition for correlated noise mechanisms. Informally, conditional composition says the mechanism's privacy guarantee is the composition of high-probability privacy guarantees for each round, which will be much stronger than worst-case privacy guarantees. 
We believe that conditional composition will be a useful framework for analyzing correlated noise mechanisms beyond those studied in this work.

\mypar{Matrix mechanism with subsampling (\cref{sec:cptb})} 
Our main contribution is to apply our conditional composition theorem to give amplified privacy guarantees for the matrix mechanism with subsampling.

We give a brief informal explanation of how this analysis proceeds:
For the analysis, instead of thinking about adding correlated noise $\bfC^{-1} \bfz$ to the subsampled data $\bfx$, it is easier to think about adding independent noise $\bfz$ to a transformation of the subsampled data $\bfC \bfx$; these views are equivalent.
To keep things simple, we can think of $\bfx \in \{0,1\}^n$, where $\bfx_i=1$ indicates that the target person's data was included in minibatch $i$. The subsampling probability is $\Pr[\bfx_i=1]=p$ independently for each $i$.
Our analysis applies conditional composition over the coordinates of $\bfy = \bfC \bfx + \bfz$. The noise $\bfz$ is independent across coordinates, but the subsampling is not. That is, conditioned on $\bfy_1, \cdots, \bfy_{i-1}$, the distribution of $\bfx$ is no longer a product distribution with $\Pr[\bfx_j=1]=p$. However, we can show that, with high probability over $\bfy_1, \cdots, \bfy_{i-1}$, the conditional distribution of $\bfx$ is well-behaved. 
Specifically, $\Pr[\bfx_j=1|\bfy_1, \cdots, \bfy_{i-1}]\le p_{i,j}$ with high probability over $\bfy_1, \cdots, \bfy_{i-1}$.
This allows us to use standard tools give a conditional amplified privacy guarantee (with $p_{i,j}$ as the subsampling probability instead of $p$), which we can then pass to our conditional composition theorem.\tnote{This paragraph still feels too technical...}

The brief explanation above ignores some technicalities:
The subsampled data $\bfx$ is not a sequence of bits, but a sequence of adaptive vectors. That is, $\bfx_i$ is actually the sum of the clipped gradients in the random minibatch at round $i$.
Nevertheless we are able to show that the general case is dominated by $\bfx$ being a sequence of bits
To do this, we show that the conditional distribution of each round of the matrix mechanism with subsampling is a \textit{mixture of Gaussians (MoG)}. We next prove some basic properties about MoG mechanisms which allow us to show that with high probability, each round of the matrix mechanism is dominated by a given a MoG mechanism. Then by the results in \cref{sec:conditioningtrick}, the privacy guarantee of the composition of these dominating MoG mechanisms is a valid privacy guarantee for the matrix mechanism. Our analysis yields an algorithm, \MMCC{} given in \cref{fig:cptb}, for computing the dominating MoG mechanisms and their composition's privacy guaranteeis. 

\MMCC{} has some nice quantitative properties we discuss here. First, it is nearly-tight in the low-epsilon regime, i.e., the limit as $\epsilon \rightarrow 0$ of $\epsilon$ computed by \MMCC{} divided by true epsilon is 1. Second, it quantitatively improves on the only other work on amplification for matrix mechanisms \citep{BandedMF} in multiple ways. Unlike the results of \citet{BandedMF}, \MMCC{} can compute a nearly-tight $\epsilon$ value for full lower-triangular matrices if $\epsilon$ is small, and in particular nearly retrieves the privacy guarantee of DP-SGD with sampling if $\bfC$ is not banded but still close to the identity. In addition, \MMCC{} analyzes i.i.d.~sampling every round, i.e., the ``most random'' form of sampling. The added randomness allows for more gains from privacy amplification than the more structured (i.e., less random) sampling scheme analyzed by \cite{BandedMF}.

\mypar{Binary tree mechanism conditional composition (\cref{sec:ftrl})} To demonstrate the versatility of conditional composition, we use it to analyze the binary tree mechanism with privacy amplification via shuffling (instead of sampling). \cite{kairouz2021practical} uses a post-processing of the unamplified binary tree mechanism for regret minimization called TreeAgg-DP-FTRL, which vastly improves on the regret guarantees of unamplified DP-SGD. However, in centralized settings, the noise added in TreeAgg-DP-FTRL is worse than amplified DP-SGD can be worse by a factor of $\sqrt{\log n}$ (for sufficiently small $\epsilon$). We nearly close this gap by giving an amplification analysis for TreeAgg-DP-FTRL. In particular, we show that TreeAgg-DP-FTRL with privacy amplification by shuffling allows us to use $\sigma$ that is smaller by a multiplicative $\sqrt{\log n} / \sqrt{\log \log (1/\delta)}$ (again for sufficiently small $\epsilon$) than in the unamplified case. This is nearly tight - \agnoteinline{@AT can you fill in the argument here?}
\anoteinline{What does it means to have an amplified algorithm have a regret guarantee? Regret is defined by streaming data where the ordering is fixed.}
\agnoteinline{Changed it to noise, but I think it's stronger if we can say some utility statement had a sqrt(log n) gap we close}
\anoteinline{We also need to mention that ${\log n}$ is the best amplification we can hope for, since for the tree, we output independent $\log n$ histogragms. Hnce the error for each of the ones cannot be better than $\Omega(\frac{\sqrt{\log(1/\delta)}}{\epsilon})$. And that is pretty much what we are getting.}

\mypar{Empirical $\epsilon$ computations (\cref{sec:empirical})} We implement \MMCC{} and use it to compute $\epsilon$ for the binary tree mechanism, as well as the matrix mechanisms studied in \cite{HU22} and \cite{BandedMF}. For the first two (non-banded) settings, we demonstrate that the improvement in $\epsilon$ from sampling computed by \MMCC{} is $\Omega(\sqrt{\log n})$, matching the theoretical predictions from \cref{sec:ftrl}. For the matrices in \cite{BandedMF}, we compare to their amplification argument and sampling scheme. While we do not improve over their amplification results for general banded matrices, empirically we show that when the columns of $\bfC$ are decreasing, \MMCC{} computes tighter epsilons than the techniques in \cite{BandedMF}.

\mypar{Learning experiments (\cref{sec:experiments})} Equipped with \MMCC{}, we revisit the experiments in \cite{BandedMF}, where the authors showed that their MF-DP-FTRL with their amplification scheme is never worse than amplified DP-SGD for training on CIFAR-10. However, for $\epsilon < 1$, amplified DP-SGD remains the best algorithm. Since \MMCC{} can handle i.i.d sampling every round, which is more random than the sampling in \cite{BandedMF}, we expect that the minimum $\epsilon$ for which MF-DP-FTRL with i.i.d sampling and $\epsilon$ computed by \MMCC{} improves over DP-SGD is smaller than in \cite{BandedMF}. We conduct experiments to verify this conjecture. \agnote{TODO: Fill in once experiments are done.}
